\DocumentMetadata{pdfversion=1.7,pdfstandard=A-3a,lang=zh-CN,tagging=on}
\documentclass[12pt,fontset=windows,UTF8]{ctexart}
\usepackage[a4paper,hmargin=1cm,vmargin=2cm]{geometry}
\usepackage[main=chinese,english]{babel}
\usepackage{amsmath}
\usepackage{unicode-math}
\usepackage{tabularx}
\usepackage{booktabs}
\usepackage{float}

\setmathfont{STIX Two Math}

\title{工程数学公式}
\author{石成玉}
\date{\today}

\begin{document}

\section{公式}

\subsection{欧拉公式}

\begin{equation}
	e^{ix}=\cos x + i \sin x \label{eq:欧拉公式}
\end{equation}

\subsection{三角函数公式}

\subsubsection{和差化积公式}

\begin{equation}
	\begin{aligned}
		\sin \alpha + \sin \beta & = 2 \sin \frac{\alpha+\beta}{2} \cos\frac{\alpha+\beta}{2} \\
		\sin \alpha - \sin \beta & = 2 \cos \frac{\alpha+\beta}{2} \sin\frac{\alpha+\beta}{2} \\
		\cos\alpha + \cos\beta & = 2\cos\frac{\alpha+\beta}{2}\cos\frac{\alpha-\beta}{2} \\
		\cos\alpha - \cos\beta & = -2\cos\frac{\alpha+\beta}{2}\cos\frac{\alpha-\beta}{2} \\
		\tan\alpha + \tan\beta & = \frac{\sin(\alpha+\beta)}{\cos\alpha\cos\beta}
	\end{aligned}
\end{equation}

\subsubsection{积化和差公式}

\begin{equation}
	\begin{aligned}
		\cos\alpha\sin\beta & = \frac{1}{2}\left[\sin(\alpha+\beta) - \sin(\alpha+\beta)\right] \\
		\sin\alpha\sin\beta & =-\frac{1}{2}\left[\cos(\alpha+\beta)-\cos(\alpha-\beta)\right] \\
		\cos\alpha\cos\beta & =\frac{1}{2}\left[\cos(\alpha+\beta)-\cos(\alpha-\beta)\right]
	\end{aligned}
\end{equation}

\subsubsection{二倍角公式}

\begin{equation}
	\begin{aligned}
		\sin(2\alpha) & =2\sin\alpha\cos\alpha=\frac{2\tan\alpha}{1+\tan^2\alpha} \\
		\cos(2\alpha) & = \cos^2\alpha - \sin^2\alpha = 2\cos^2\alpha - 1 = 1 - 2\sin^2\alpha = \frac{1-\tan^2\alpha}{1+\tan^2\alpha} \\
		\tan(2\alpha) & = \frac{2\tan\alpha}{1-\tan^2\alpha}
	\end{aligned}
\end{equation}

\section{分类问题}

\begin{table}[H]
	\caption{定解问题}
	\label{tb:定解问题}
	\begin{tabularx}{\linewidth}{
		>{\hsize=0.4\hsize\linewidth=\hsize\centering\arraybackslash}X
		>{\hsize=0.2\hsize\linewidth=\hsize\centering\arraybackslash}X
		>{\hsize=0.2\hsize\linewidth=\hsize\centering\arraybackslash}X
		>{\hsize=0.2\hsize\linewidth=\hsize\centering\arraybackslash}X
	}
		\toprule
		方程 & 边界条件 & 初值条件 & 方程类型 \\
		\midrule
		波动方程 & 有 & 有 & 双曲型方程 \\
		热传导方程 & 有 & 有 & 抛物型方程 \\
		拉普拉斯方程 & 有 & 无 & 椭圆型方程 \\
		泊松方程 & 有 & 无 & 椭圆型方程 \\
		\bottomrule
	\end{tabularx}
\end{table}

\end{document}